\documentclass[11pt,a4paper]{article}

\usepackage[utf8]{inputenc}
\usepackage[T1]{fontenc}
\usepackage[spanish,es-noshorthands]{babel}
\usepackage{xcolor}
\usepackage[most]{tcolorbox}
\usepackage{listings}
\usepackage{fancyhdr}
\usepackage{lastpage}
\usepackage{enumitem}
\usepackage[hidelinks]{hyperref}

\setlength{\headheight}{30pt}
\pagestyle{fancy}
\fancyhf{}
\renewcommand{\headrulewidth}{0pt}
\fancyhead[R]{%
  \fontsize{8}{9.6}\selectfont
  Licenciatura en Ciencias de la Computación\\
  Sistemas Operativos\\[2pt]
  \rule{4.2cm}{0.5pt}
}
\fancyfoot[C]{\thepage\ / \pageref{LastPage}}

\newtcolorbox{trackbox}[1]{
  colback=green!7!white, colframe=green!45!black,
  arc=2mm, boxrule=0.6pt, left=8pt, right=8pt, top=6pt, bottom=6pt,
  title=TRACK -- #1, fonttitle=\bfseries, coltitle=green!35!black,
  colbacktitle=green!7!white, boxsep=1pt, breakable
}

\newtcolorbox{extrabox}{
  colback=orange!8!white, colframe=orange!55!white,
  arc=2mm, boxrule=0.6pt, left=8pt, right=8pt, top=6pt, bottom=6pt,
  title=EXTRA, fonttitle=\bfseries, coltitle=orange!55!black,
  colbacktitle=orange!8!white, boxsep=1pt, breakable
}

\newtcolorbox{pointsextrabox}{
  colback=red!7!white, colframe=red!55!white,
  arc=2mm, boxrule=0.6pt, left=8pt, right=8pt, top=6pt, bottom=6pt,
  title=BONUS · POR PUNTOS \coin, fonttitle=\bfseries, coltitle=red!55!black,
  colbacktitle=red!7!white, boxsep=1pt, breakable
}

\newcommand{\coin}{\raisebox{0.1ex}{\tikz{\draw[fill=yellow!60!orange] (0,0) circle (0.5ex);}}}

\lstdefinestyle{codestyle}{
  backgroundcolor=\color{black!5}, rulecolor=\color{black!25},
  frame=single, framerule=0.5pt, basicstyle=\ttfamily\small,
  breaklines=true, showstringspaces=false, tabsize=4,
  columns=fullflexible
}
\lstset{style=codestyle}

\title{Trabajo Práctico 2 -- Ejercicios extra}
\author{}
\date{}

\begin{document}
\renewcommand{\contentsname}{Tracks}
\maketitle
\thispagestyle{fancy}
\begin{center}
SISTEMAS OPERATIVOS
\end{center}
\newpage

\tableofcontents
\newpage

Estos tracks amplían la práctica de procesos con problemas pequeños pero cercanos a herramientas reales. Podés empezar por cualquier track. Cada uno tiene una parte principal y una extensión para seguir investigando.

\section{Señales y control de trabajos (5/10)}

\begin{trackbox}{5/10}
\textbf{Descripción:} usar señales para pausar, continuar y terminar procesos. Es una base práctica para administrar servidores y depurar programas.
\end{trackbox}

\begin{enumerate}
  \item Ejecutar \texttt{sleep 100} en foreground y detenerlo con \texttt{Ctrl-Z}.
  \item Consultar el trabajo con \texttt{jobs}, continuarlo en background con \texttt{bg} y traerlo nuevamente con \texttt{fg}.
  \item Encontrar su PID y enviar \texttt{SIGSTOP}, \texttt{SIGCONT} y \texttt{SIGTERM} usando \texttt{kill}.
  \item Comparar \texttt{SIGTERM} con \texttt{SIGKILL}. Explicar por qué un programa no puede capturar \texttt{SIGKILL}.
  \item Escribir un programa que instale un handler para \texttt{SIGINT}, cuente cuántas veces recibió la señal y termine limpiamente con la segunda.
\end{enumerate}

\begin{extrabox}
Investigá la diferencia entre una señal estándar y una señal encolada. Usá \texttt{kill -l}, \texttt{man 7 signal} y \texttt{sigaction()}.
\end{extrabox}

\newpage
\section{Mini-shell con \texttt{fork()} y \texttt{exec()} (7/10)}

\begin{trackbox}{7/10}
\textbf{Descripción:} construir una versión mínima del mecanismo que usa un shell real. Integra parsing, creación de procesos, ejecución y espera.
\end{trackbox}

\begin{enumerate}
  \item Leer una línea usando \texttt{fgets()} y separar el comando de sus argumentos.
  \item Crear un hijo con \texttt{fork()} y ejecutar el comando con \texttt{execvp()}.
  \item Hacer que el padre espere al hijo con \texttt{waitpid()}.
  \item Implementar el comando interno \texttt{cd}, que debe ejecutarse en el proceso del shell.
  \item Agregar ejecución en background cuando la línea termina en \texttt{\&}.
  \item Informar errores de \texttt{fork()}, \texttt{execvp()} y \texttt{waitpid()} con \texttt{perror()}.
\end{enumerate}

\begin{lstlisting}[language=C]
if (fork() == 0) {
    execvp(args[0], args);
    perror("execvp");
    _exit(127);
}
wait(NULL);
\end{lstlisting}

\begin{pointsextrabox}
Agregá una tubería entre dos comandos usando \texttt{pipe()}, \texttt{dup2()} y dos procesos hijo. Probá ejecutar \texttt{ps aux | grep bash}.
\end{pointsextrabox}

\newpage
\section{Carreras y sincronización entre hilos (7/10)}

\begin{trackbox}{7/10}
\textbf{Descripción:} observar una condición de carrera y corregirla con un mutex. Es una práctica directamente útil para programación concurrente.
\end{trackbox}

\begin{enumerate}
  \item Crear diez hilos que incrementen un contador global un millón de veces cada uno.
  \item Ejecutar varias veces y comparar el resultado con el valor esperado.
  \item Explicar por qué \texttt{contador++} no es una operación indivisible.
  \item Proteger el incremento usando un mutex:
    \begin{lstlisting}[language=C]
pthread_mutex_lock(&mutex);
contador++;
pthread_mutex_unlock(&mutex);
    \end{lstlisting}
  \item Medir el tiempo con y sin mutex usando \texttt{time}.
  \item Confirmar la cantidad de hilos con \texttt{ps -T}.
\end{enumerate}

\begin{extrabox}
Compará el mutex con operaciones atómicas usando \texttt{stdatomic.h}. Investigá por qué una solución correcta puede ser más lenta y cuándo conviene cada alternativa.
\end{extrabox}

\newpage
\section{Cambios de contexto y rendimiento (8/10)}

\begin{trackbox}{8/10}
\textbf{Descripción:} medir qué costo tiene alternar entre procesos e hilos. Conecta el modelo teórico del planificador con datos observables.
\end{trackbox}

\begin{enumerate}
  \item Crear un programa que haga trabajo de CPU durante varios segundos.
  \item Medir cambios de contexto voluntarios y no voluntarios antes y después usando \texttt{/proc/PID/status}.
  \item Comparar un proceso que usa \texttt{sleep()} con otro que ocupa la CPU en un loop.
  \item Ejecutar dos copias del programa y observar el efecto sobre cambios de contexto y tiempo de ejecución.
  \item Repetir la medición con dos hilos y explicar qué resultados esperás encontrar.
  \item Registrar resultados en una tabla y distinguir una hipótesis de una observación.
\end{enumerate}

\begin{pointsextrabox}
Usá \texttt{taskset} para fijar la afinidad a un núcleo y repetí el experimento. Investigá también \texttt{getrusage()} y la diferencia entre tiempo de usuario y tiempo de sistema.
\end{pointsextrabox}

\end{document}
